% !TeX root = ../navier-stokes-manuscript.tex
% ============================================================
% 1.3 Gold Standard Proof
% ============================================================
\subsection{Gold Standard Proof}\label{sub:GoldStandardProof}

The historical results leave a clean direct route to global smoothness.
Let \([0,T_*)\) be the maximal interval of the smooth solution generated by \(u_0\).
A Gold-standard proof derives, for some fixed \(s_0>5/2\),
\[
  \sup_{0\leq t<T_*}
  \lVert u(t)\rVert_{H^{s_0}(\Omega)}
  <
  \infty
\]
from \(u_0\), \(\nu\), and the Navier--Stokes equations themselves.
This is the direct estimate absent from the classical Gold route surveyed above.

The payoff is larger than the single norm might suggest.
Sobolev embedding gives \(H^{s_0}\hookrightarrow W^{1,\infty}\), so the bound controls the velocity gradient that appears in every higher differentiated transport estimate.
For each fixed finite \(m\geq s_0\),
\[
  \frac12\frac{d}{dt}\lVert u\rVert_{H^m}^{2}
  +\nu\lVert\nabla u\rVert_{H^m}^{2}
  \leq
  C_m\lVert\nabla u\rVert_{L^\infty}\lVert u\rVert_{H^m}^{2}.
\]
The initial datum already belongs to every \(H^m\), and Gronwall's inequality propagates each finite order to \(T_*\).
The constants may depend on \(m\), because \(C^\infty\) requires every finite derivative and imposes no common numerical bound on the entire infinite family.
Higher Sobolev embeddings then recover all continuous spatial derivatives.

Pressure is recovered from the same velocity through
\[
  -\Delta p
  =
  \sum_{i,j=1}^{3}\partial_i\partial_j(u_i u_j),
\]
and the momentum equation recovers \(\partial_tu\).
Repeated differentiation recovers every higher mixed derivative of \(u\) and \(p\).
A late time slice can then be used as new initial data for local theory, and the uniform \(H^{s_0}\) bound gives enough lifespan to cross \(T_*\).
This contradicts maximality, so the direct estimate forces \(T_*=\infty\).

The obstruction appears as soon as one tries to produce that estimate from the basic energy law.
Energy gives
\[
  \lVert u(t)\rVert_{L^2}^{2}
  +2\nu\int_0^t\lVert\nabla u(s)\rVert_{L^2}^{2}\,ds
  =
  \lVert u_0\rVert_{L^2}^{2}.
\]
It controls the total time integral of the squared gradient.
Raising the calculation by one spatial derivative gives
\[
  \frac12\frac{d}{dt}\lVert\nabla u\rVert_{L^2}^{2}
  +\nu\lVert\Delta u\rVert_{L^2}^{2}
  =
  \int_\Omega (u\cdot\nabla)u\cdot\Delta u\,dx.
\]
The transport term now survives, and the usual interpolation and Young inequalities give
\[
  \left|
    \int_\Omega (u\cdot\nabla)u\cdot\Delta u\,dx
  \right|
  \leq
  \frac{\nu}{2}\lVert\Delta u\rVert_{L^2}^{2}
  +C\nu^{-3}\lVert\nabla u\rVert_{L^2}^{6}.
\]
After viscosity absorbs the first term, the estimate asks for time control of the sixth power of the gradient.
The energy identity supplies time control only of its second power.
Tall, narrow peaks can respect the second-power integral while defeating the sixth-power integral, so this estimate does not close the first derivative level.
That observation identifies an estimate gap; it does not assert that the fluid actually forms those peaks.

The Gold burden is easy to state and hard to pay.
One must derive a continuation-grade bound from the coupled action of transport, pressure, incompressibility, and viscosity without assuming the strong quantity the estimate is meant to control.
Once that payment is made, the Sobolev bootstrap, pressure recovery, and local restart complete the argument.
Silver keeps the same smoothness claim and changes the direction of attack by asking what any alleged nonsmooth outcome must have ceased to satisfy.
